\chapter{Conformal geometry of the Willmore energy}\label{sec:willmore-conformal-calculus}
This section expands Definition \ref{def:willmore-energy}. All integrals are over the parameter surface with its induced metric, retaining immersion multiplicity. We use the trace mean curvature vector $\vec H=\Delta F$ in Euclidean space. For a two-sided hypersurface, $A=D\nu$, $H=\operatorname{tr}A$ and $\vec H=-H\nu$. The conformal mechanism is described in \cite{LY82}, Proposition 1, p.~279, and the spherical formulation in \cite{MNWill}, pp.~684--685. The formulas below are derived explicitly.

\begin{definition}[Vector normalization and the trace-free second fundamental form]\label{def:willmore-vector-normalization}
For a smooth immersion $F:M^2\to\mathcal N^d$ into a space of constant sectional curvature $\kappa$, let $B$ be its normal-vector-valued second fundamental form, $\vec H=\operatorname{tr}_gB$, and
\[
 B^0=B-\tfrac12g\otimes\vec H,
 \qquad \mathcal W_\kappa(F)=\int_M\left(\kappa+\tfrac14|\vec H|^2\right)d\mu.
\]
For a closed Euclidean surface write $\mathcal W=\mathcal W_0$. In a two-sided three-dimensional ambient space, $B=-A\nu$, so $|\vec H|^2=H^2$. Li--Yau use averaged mean curvature in \cite{LY82}; Simon uses the trace vector and the factor $1/4$ in \cite{SimW93}, p.~281.
\end{definition}

\begin{lemma}[Conformal changes of the normal geometry]\label{lem:conformal-second-fundamental-form}
Let $\widetilde{\bar g}=e^{2\phi}\bar g$ be a smooth conformal change of ambient metric near an immersed surface, and set $\varphi=\phi\circ F$. Tangent and normal spaces as vector subspaces are unchanged. Regarding normal-valued tensors as vectors in the same ambient tangent bundle,
\[
 \widetilde g=e^{2\varphi}g,\qquad
 \widetilde B(U,V)=B(U,V)-g(U,V)(\bar\nabla\phi)^\perp,
\]
\[
 \widetilde{\vec H}=e^{-2\varphi}(\vec H-2(\bar\nabla\phi)^\perp),
 \qquad \widetilde B^0=B^0,
 \qquad |\widetilde B^0|_{\widetilde g,\widetilde{\bar g}}^2\,d\widetilde\mu=|B^0|_{g,\bar g}^2\,d\mu.
\]
For a hypersurface with $\widetilde\nu=e^{-\phi}\nu$, the shape operators instead satisfy $\widetilde A=e^{-\phi}(A+\nu(\phi)I)$.
\end{lemma}
\begin{proof}
The ambient connection changes by
\[
 \widetilde{\bar\nabla}_UV=\bar\nabla_UV+U(\phi)V+V(\phi)U-\bar g(U,V)\bar\nabla\phi.
\]
To check this identity, the right-hand connection is torsion free, and substituting it into the derivative of $e^{2\phi}\bar g(V,W)$ gives metric compatibility. Uniqueness of the Levi-Civita connection proves the formula. The two middle terms are tangent when $U,V$ are tangent to the surface. Taking the normal part gives the formula for $\widetilde B$. Tracing uses $\widetilde g^{-1}=e^{-2\varphi}g^{-1}$, which gives the stated vector mean curvature. Subtracting $\widetilde g\otimes\widetilde{\vec H}/2$ cancels the added pure-trace normal term.

For the norm of a normal-valued bilinear form, two inverse induced metrics contribute $e^{-4\varphi}$ and the ambient vector norm contributes $e^{2\varphi}$. Hence the squared norm of $B^0$ changes by $e^{-2\varphi}$; the area element changes by $e^{2\varphi}$. Finally differentiate $e^{-\phi}\nu$ using the connection formula and take its tangential part. The terms involving $U(\phi)\nu$ cancel, leaving $e^{-\phi}(AU+\nu(\phi)U)$, as claimed. This last scalar formula uses the convention $A=D\nu$.
\end{proof}

\begin{lemma}[The intrinsic curvature density under a conformal change]\label{lem:surface-conformal-gauss-density}
For a smooth surface metric $\widetilde g=e^{2\varphi}g$,
\[
 \widetilde K=e^{-2\varphi}(K-\Delta_g\varphi),\qquad
 \widetilde K\,d\widetilde\mu=(K-\Delta_g\varphi)\,d\mu.
\]
Consequently the integrals of Gauss curvature agree on a closed surface.
\end{lemma}
\begin{proof}
The difference of induced connections is
$C_{ij}^{k}=\varphi_i\delta_j^k+\varphi_j\delta_i^k-g_{ij}\varphi^k$.
At a point choose $g$-geodesic coordinates with orthonormal coordinate vectors. The three relevant coefficient vectors are
\[
 (C_{11}^1,C_{11}^2)=(\varphi_1,-\varphi_2),\quad
 (C_{12}^1,C_{12}^2)=(\varphi_2,\varphi_1),\quad
 (C_{22}^1,C_{22}^2)=(-\varphi_1,\varphi_2).
\]
In the component of $\widetilde R(e_1,e_2)e_2$ along $e_1$, the derivative terms contribute $-\varphi_{11}-\varphi_{22}$. The products of the displayed $C$ coefficients cancel. The remaining curvature term is $K$. Pairing in $\widetilde g$ contributes $e^{2\varphi}$ and dividing by the squared area of the coordinate parallelogram contributes $e^{-4\varphi}$. This gives the first identity and then the density identity. On a closed surface the integral of $\Delta_g\varphi$ vanishes by the divergence theorem, including on a nonorientable surface using its Riemannian measure or orientation double cover.
\end{proof}

\begin{proposition}[The conformally invariant energy density modulo curvature]\label{prop:willmore-tracefree-decomposition}
For a surface in a space of constant sectional curvature $\kappa$,
\[
 \kappa+\tfrac14|\vec H|^2=K+\tfrac12|B^0|^2.
\]
If an ambient conformal diffeomorphism maps a neighborhood of a closed immersed surface in one space form into another, its Willmore energy with the corresponding ambient curvature term is unchanged.
\end{proposition}
\begin{proof}
In an orthonormal tangent frame the Gauss equation reads
\[
 K=\kappa+\langle B_{11},B_{22}\rangle-|B_{12}|^2
 =\kappa+\tfrac12(|\vec H|^2-|B|^2).
\]
The orthogonal trace decomposition gives $|B|^2=|B^0|^2+|\vec H|^2/2$, proving the first assertion. Pull back the target ambient metric by the conformal map. Lemma \ref{lem:conformal-second-fundamental-form} preserves the trace-free norm density, and Lemma \ref{lem:surface-conformal-gauss-density} preserves the integral of Gauss curvature. Integrate the displayed identity on each side. The ambient map must be smooth on a neighborhood of the entire image; a pole lying on the surface is excluded here and is treated separately below. This proof does not require the Gauss--Bonnet theorem.
\end{proof}

\begin{proposition}[Stereographic projection preserves the normalized energy]\label{prop:stereographic-willmore-identity}
Let $F:M^2\to S^N\setminus\{p\}$ be a smooth closed immersion in the unit sphere, and let $\pi_p$ be stereographic projection. Then
\[
 \mathcal W_0(\pi_p\circ F)=\mathcal W_1(F)
 =\int_M\left(1+\tfrac14|\vec H_{S^N}|^2\right)d\mu.
\]
In particular, the Euclidean energy of the projection of a closed minimal surface in the sphere equals its spherical area. The identity is equation (1) of \cite{MNWill}, p.~684, in trace-vector notation.
\end{proposition}
\begin{proof}
Take $p=(0,\ldots,0,1)$. The inverse projection is
\[
 Y(x)=\left(\frac{2x}{1+|x|^2},\frac{|x|^2-1}{1+|x|^2}\right).
\]
A direct derivative gives, for every $v$,
\[
 |dY_x(v)|^2=\frac{4}{(1+|x|^2)^2}|v|^2.
\]
For example, the first component has derivative $2v/(1+r^2)-4x(x\cdot v)/(1+r^2)^2$ and the last has derivative $4(x\cdot v)/(1+r^2)^2$; their squared cross terms cancel. Thus projection is conformal. Compactness and avoidance of $p$ ensure a smooth conformal factor near the entire surface, so Proposition \ref{prop:willmore-tracefree-decomposition} applies with ambient curvatures one and zero. Minimality removes the mean-curvature term in the spherical integral.
\end{proof}

\begin{theorem}[First variation of Willmore energy in a three-dimensional space form]\label{thm:willmore-first-variation-spaceform}
Let $F:M^2\to N^3_\kappa$ be a smooth closed two-sided immersion and let its normal variation have initial velocity $f\nu$. With $H=\operatorname{tr}A$ and $A^0=A-HI/2$,
\[
 \left.\frac{d}{dt}\right|_{0}\mathcal W_\kappa(F_t)
 =-\frac12\int_M\bigl(\Delta H+H|A^0|^2\bigr)f\,d\mu.
\]
Thus critical points satisfy $\Delta H+H|A^0|^2=0$. In averaged notation $h=H/2$, this is
\[
 \Delta h+2h(h^2+\kappa-K)=0.
\]
This is the Willmore variational problem discussed in \cite{MNWill}, p.~685; the displayed formula retains the ambient curvature term explicitly.
\end{theorem}
\begin{proof}
The area variation is $\dot{d\mu}=Hf\,d\mu$. The mean-curvature linearization is $\dot H=-\Delta f-(|A|^2+2\kappa)f$, since the ambient Ricci curvature in the normal direction is $2\kappa$. Differentiate the energy to obtain
\[
 \dot{\mathcal W}_\kappa
 =\int\left[-\tfrac12H\Delta f-\tfrac12H(|A|^2+2\kappa)f
 +(\kappa+H^2/4)Hf\right]d\mu.
\]
The two explicit $\kappa Hf$ terms cancel. Self-adjointness of the Laplacian on a closed surface moves $\Delta$ from $f$ to $H$. Since $|A|^2-H^2/2=|A^0|^2$, this is the asserted formula. The Gauss equation $K=\kappa+\det A$ gives $|A^0|^2=2(h^2+\kappa-K)$, proving the averaged expression. Arbitrary smooth normal speeds imply the Euler--Lagrange equation. Tangential velocities only reparametrize the immersion and do not change the energy.
\end{proof}

\begin{corollary}[Minimal spherical surfaces give Euclidean Willmore surfaces]\label{cor:minimal-spherical-willmore}
A closed two-sided minimal surface in a three-dimensional space form is a Willmore critical point. A stereographic projection of a closed minimal surface of $S^3$ avoiding the pole is a Euclidean Willmore critical point. Its energy equals the original area.
\end{corollary}
\begin{proof}
The equation of Theorem \ref{thm:willmore-first-variation-spaceform} vanishes when $H=0$. For the second assertion, every sufficiently small smooth variation of the projected compact surface pulls back to a spherical variation avoiding the pole. Proposition \ref{prop:stereographic-willmore-identity} identifies their energies for all small parameters. Differentiating therefore transfers stationarity. If the spherical immersion is one-sided, apply the preceding first-variation formula on its normal orientation double cover: the lifted energy and its first derivative are twice those downstairs, so stationarity descends. The energy value follows from the same identity. This only proves criticality, not minimizing or stability among all variations.
\end{proof}
